


\section{Conclusion}
% This work proposes Chinese CLIP, an implementation of CLIP~\citep{clip} pretrained on large-scale Chinese datasets with contrastive learning. 
In this work, we propose Chinese CLIP, a Chinese-specific vision-language foundation model. 
Specifically, we construct a pretraining dataset of around $200$ million samples, and pretrain a series of Chinese CLIP models with the proposed two-stage pretraining method, which improves both pretraining efficiency and effectiveness. 
% We implement pretraining by initializing encoders with the weights of the existing pretrained models and using our proposed two-stage pretraining method to build $5$ Chinese CLIP models of different model sizes, ranging from $77$ to $958$ million parameters. 
Our comprehensive evaluation shows that Chinese CLIP can reach state-of-the-art performance on multiple cross-modal retrieval datasets in zero-shot learning and finetuning. 
Furthermore, we demonstrate that Chinese CLIP models can also achieve competitive performance in zero-shot image classification across $10$ datasets. 

% In the future, we will explore more efficient methods for pretraining and finetuning Chinese CLIP, and we plan to build vision-language understanding benchmarks based on Chinese-native data, which can better reflect the effectiveness of language-specific vision-language models. 

% However, we still find that Chinese CLIP is limited in some aspects, such as sensitivity to prompts and difficulty in understanding negation. 
% Besides, there is still much room in exploring data scaling and model scaling. 
% In the future, we will continue our research in building Chinese CLIP models with better performance and exploring the features of contrastive-learning-based pretrained models. 